% Part: applied-modal-logics
% Chapter: epistemic-logic
% Section: language-epistemic-logic

\documentclass[../../../include/open-logic-section]{subfiles}

\begin{document}

\olfileid{aml}{el}{lan}

\olsection{认知逻辑的语言}

\begin{defn}
令 $G$ 为一个主体符号集。多主体认知逻辑的基本语言包含：
\begin{enumerate}
  \tagitem{prvFalse}{表示假的命题常项~$\lfalse$。}{}
  \tagitem{prvTrue}{表示真的命题常项~$\ltrue$。}{}
  \item 一个可数无穷的命题变元集：$\Obj p_0$, $\Obj p_1$, $\Obj p_2$, \dots
  \item 命题联结词：\startycommalist
  \iftag{prvNot}{\ycomma $\lnot$（否定）}{}%
  \iftag{prvAnd}{\ycomma $\land$（合取）}{}%
  \iftag{prvOr}{\ycomma $\lor$（析取）}{}%
  \iftag{prvIf}{\ycomma $\lif$（条件句）}{}%
  \iftag{prvIff}{\ycomma $\liff$（双条件句）}{}.
  \item 知识算子 $\Knows_a$，其中 $a \in G$。
\end{enumerate}
\end{defn}

如果我们只关心系统中单个主体的知识，就可以省去对集合~$G$ 及各个主体的提及。在这种情况下，我们只需使用基本算子~$\Knows$。

\begin{defn}
认知语言的\emph{公式}归纳定义如下：
\begin{enumerate}
\tagitem{prvFalse}{$\lfalse$ 是原子公式。}{}

\tagitem{prvTrue}{$\ltrue$ 是原子公式。}{}

\item 每个命题变元 $\Obj p_i$ 都是（原子）公式。

\tagitem{prvNot}{如果 $!A$ 是公式，则 $\lnot !A$ 是公式。}{}

\tagitem{prvAnd}{如果 $!A$ 和 $!B$ 是公式，则 $(!A \land
  !B)$ 是公式。}{}

\tagitem{prvOr}{如果 $!A$ 和 $!B$ 是公式，则 $(!A \lor !B)$
  是公式。}{}

\tagitem{prvIf}{如果 $!A$ 和 $!B$ 是公式，则 $(!A \lif !B)$
  是公式。}{}

\tagitem{prvIff}{如果 $!A$ 和 $!B$ 是公式，则 $(!A \liff !B)$
  是公式。}{}

\item 如果 $!A$ 是公式且 $a \in G$，则 $\Knows_a !A$ 是公式。

\tagitem{limitClause}{其他都不是公式。}{}
\end{enumerate}
\end{defn}

如果一个公式~$!A$ 不包含 $\Knows_a$，则称它是\emph{不含模态词的}。

\begin{defn}
  $\Knows$ 算子用于表示个体知识，而 $\EKnows$（通常读作“每个人都知道”）则表示群体知识。当 $G' \subseteq G$ 时，我们将 $\EKnows_{G'} !A$ 定义为 \[\bigwedge_{b \in G'} \Knows_b !A.\] 的缩写。
\end{defn}

我们还可以定义一种更强的知识概念，即一组主体~$G$ 之间的\emph{公共知识}。当一条信息在某个主体群体中是公共知识时，意味着对该群体中主体的任意组合，所有人都知道所有人都知道……以至无穷。这比群体知识要强得多，并且很容易构造出某个公式是群体知识而非公共知识的关系模型。我们将用 $\CKnows_G !A$ 来表示“$!A$ 在~$G$ 中是公共知识”。

\end{document}